% Source PDF page 36; manual page 2-7.
\subsection{Earth-Centered, Inertial, Cartesian (ECIC) Coordinate System}\label{sec:ecic}

The center of the earth is assumed to be nonaccelerating, so an inertial coordinate
system can be defined with the earth's center as its origin. The inertial $z$ axis,
$\hat{z}_{I}$, is aligned with the earth's spin axis $\hat{z}_{\oplus}$ and is
therefore fixed with respect to the earth.

The $x$ and $y$ axes of the ECIC and ECFC systems are related by a rotation about
the common $z$ axis through the angle
\begin{equation}
  \Omega = \Omega_{0}
    + \left\lVert {}_{I}\vec{\omega}_{\oplus} \right\rVert (t-t_{0}),
  \label{eq:ecic-rotation-angle}
\end{equation}
where
\begin{equation}
  \left\lVert {}_{I}\vec{\omega}_{\oplus} \right\rVert
    = {}_{I}\vec{\omega}_{\oplus}\cdot\hat{z}_{\oplus}
    = \text{earth's rotation rate}.
  \label{eq:earth-rotation-rate}
\end{equation}
The spin axis is aligned with the ECIC and ECFC $z$ axes. The quantity $\Omega_0$
is the initial rotation angle between the two systems at time $t_0$. The initial angle,
time, and earth rotation rate are input values described in
\cref{sec:block-initial,sec:block-earth}.

Output variables giving position, velocity, and acceleration components in ECIC
coordinates are computed by the function \texttt{ecic\_coords}.

\taossubhead{Position}
\begin{equation}
\begin{aligned}
  \statevar{xecic}\quad
  \vec{r}\cdot\hat{x}_{I}
    &= \vec{r}\cdot
       (\hat{x}_{\oplus}\cos\Omega-\hat{y}_{\oplus}\sin\Omega),\\
  \statevar{yecic}\quad
  \vec{r}\cdot\hat{y}_{I}
    &= \vec{r}\cdot
       (\hat{x}_{\oplus}\sin\Omega+\hat{y}_{\oplus}\cos\Omega),\\
  \statevar{zecic}\quad
  \vec{r}\cdot\hat{z}_{I}
    &= \vec{r}\cdot\hat{z}_{\oplus}.
\end{aligned}
\label{eq:ecic-position-components}
\end{equation}

\taossubhead{Velocity}
\begin{equation}
\begin{aligned}
  \statevar{xecicdt}\quad
  \vec{V}_{I}\cdot\hat{x}_{I}
    &= \vec{V}_{\oplus}\cdot
       (\hat{x}_{\oplus}\cos\Omega-\hat{y}_{\oplus}\sin\Omega) \\
    &\quad - \left\lVert{}_{I}\vec{\omega}_{\oplus}\right\rVert
       \vec{r}\cdot
       (\hat{x}_{\oplus}\sin\Omega+\hat{y}_{\oplus}\cos\Omega),\\[0.4em]
  \statevar{yecicdt}\quad
  \vec{V}_{I}\cdot\hat{y}_{I}
    &= \vec{V}_{\oplus}\cdot
       (\hat{x}_{\oplus}\sin\Omega+\hat{y}_{\oplus}\cos\Omega) \\
    &\quad + \left\lVert{}_{I}\vec{\omega}_{\oplus}\right\rVert
       \vec{r}\cdot
       (\hat{x}_{\oplus}\cos\Omega-\hat{y}_{\oplus}\sin\Omega),\\[0.4em]
  \statevar{zecicdt}\quad
  \vec{V}_{I}\cdot\hat{z}_{I}
    &= \vec{V}_{\oplus}\cdot\hat{z}_{\oplus}.
\end{aligned}
\label{eq:ecic-velocity-components}
\end{equation}

\taossubhead{Acceleration}
\begin{equation}
\begin{aligned}
  \statevar{xecicdt2}\quad
  \vec{a}_{I}\cdot\hat{x}_{I}
    &= \frac{\vec{F}\cdot
       (\hat{x}_{\oplus}\cos\Omega-\hat{y}_{\oplus}\sin\Omega)}{m},\\
  \statevar{yecicdt2}\quad
  \vec{a}_{I}\cdot\hat{y}_{I}
    &= \frac{\vec{F}\cdot
       (\hat{x}_{\oplus}\sin\Omega+\hat{y}_{\oplus}\cos\Omega)}{m},\\
  \statevar{zecicdt2}\quad
  \vec{a}_{I}\cdot\hat{z}_{I}
    &= \frac{\vec{F}\cdot\hat{z}_{\oplus}}{m}.
\end{aligned}
\label{eq:ecic-acceleration-components}
\end{equation}
